%!TEX root = ./main.tex

\section[Tradeoff]{The Tradeoff That Does Not Go Away}

% SECTION TIMING: 33:00--41:00 (3 frames).
% TEACHING NOTE: the two tables are one table in the source talk. They are split here
% because ten rows at a readable size do not fit on one frame. Read down a column, then
% across one row. Do not read all twenty cells aloud.

\begin{frame}{They want opposite things (1/2)}
  \vspace{0.2em}
  \scriptsize
  \setlength{\tabcolsep}{4pt}
  \renewcommand{\arraystretch}{1.35}
  \begin{tabularx}{\textwidth}{>{\bfseries\raggedright\arraybackslash}p{1.55cm}
      >{\raggedright\arraybackslash}X >{\raggedright\arraybackslash}X}
    \toprule
    & \textcolor{coreorange}{\textbf{Sensing wants}} & \textcolor{ranblue}{\textbf{Communication wants}} \\
    \midrule
    Spectrum   & Wide bandwidth, for resolution & Efficient bandwidth, for rate \\
    Signal     & Deterministic waveforms & Random, information-bearing symbols \\
    Power      & High transmit power on the target & Power efficiency per bit \\
    Time       & Continuous illumination & Low-latency bursts \\
    Space      & Wide or scanning beams & Narrow beams, one per user \\
    \bottomrule
  \end{tabularx}

  \vspace{0.6em}
  \takeaway{Randomness carries information and ruins an autocorrelation.\\That single sentence is the whole conflict.}
\end{frame}

\begin{frame}{They want opposite things (2/2)}
  \vspace{0.2em}
  \scriptsize
  \setlength{\tabcolsep}{4pt}
  \renewcommand{\arraystretch}{1.35}
  \begin{tabularx}{\textwidth}{>{\bfseries\raggedright\arraybackslash}p{1.55cm}
      >{\raggedright\arraybackslash}X >{\raggedright\arraybackslash}X}
    \toprule
    & \textcolor{coreorange}{\textbf{Sensing wants}} & \textcolor{ranblue}{\textbf{Communication wants}} \\
    \midrule
    Interference & A clean echo, no clutter & A clean constellation, no radar leakage \\
    Domain     & Delay--Doppler resolution & Time--frequency multiplexing \\
    Hardware   & Precise timing, high-power amplifiers & Spectral efficiency, many users \\
    Latency    & Low-latency tracking feedback & Tolerates delay off the critical path \\
    Regulation & Frequency allocation rules & Emission and efficiency standards \\
    \bottomrule
  \end{tabularx}

  \vspace{0.6em}
  \qa{Is there a setting where both are optimal?}
     {No. There is a \textbf{Pareto frontier}: past a point, every improvement in
      one is paid for out of the other. Design work is choosing where to sit on
      it, not escaping it.}
\end{frame}

% ACCURACY NOTE: lower CRLB = better sensing, so the good corner of this plot is
% up and to the LEFT. Students read every curve as "up and to the right is better"
% and will get this backwards unless you say it.
\begin{frame}{The frontier, drawn once}
  \vspace{0.3em}
  \centering
  \begin{tikzpicture}[scale=0.95,transform shape]
    \draw[-{Latex[length=2.2mm]},softgray] (0,0) -- (6.6,0)
      node[anchor=north east,font=\scriptsize] {CRLB (sensing error floor) $\rightarrow$ worse};
    \draw[-{Latex[length=2.2mm]},softgray] (0,0) -- (0,3.6)
      node[anchor=south,font=\scriptsize] {achievable rate};
    \draw[very thick,draw=ranblue] (0.7,3.1) .. controls (2.2,2.9) and (3.0,1.9) .. (5.6,1.3);
    \node[font=\scriptsize,text=ranblue,anchor=north east] at (2.75,2.15) {attainable frontier};
    \fill[softgreen] (1.3,3.02) circle (2.6pt);
    \node[font=\scriptsize,text=softgreen,anchor=west] at (1.5,3.15) {sensing-first};
    \fill[coreorange] (4.7,1.52) circle (2.6pt);
    \node[font=\scriptsize,text=coreorange,anchor=west] at (4.9,1.75) {comms-first};
    \node[font=\scriptsize,text=softgray] at (1.8,0.7) {unreachable};
    \node[font=\scriptsize,text=softgray] at (5.5,3.0) {infeasible};
  \end{tikzpicture}

  \vspace{0.6em}
  \glossline{CRLB = Cram\'er--Rao Lower Bound: the smallest variance any unbiased
    estimator can achieve. \textbf{Lower is better}, so the good corner is up and to the left.}

  \vspace{0.2em}
  \takeaway{Every result in the rest of this lecture\\is a claim about one point on a curve like this.}
\end{frame}
